\documentclass[11pt]{article}
\usepackage[margin=1in]{geometry}
\usepackage{amsmath,amssymb,amsthm}
\usepackage{mathtools}
\usepackage{enumitem}
\usepackage{hyperref}
\usepackage{cleveref}
\usepackage{booktabs}
\usepackage{tabularx}
\usepackage{microtype}
\usepackage{cite}
\usepackage{xurl}
\usepackage{path}
\hypersetup{colorlinks=true,linkcolor=blue,citecolor=blue,urlcolor=blue}
\newtheorem{definition}{Definition}
\newtheorem{lemma}{Lemma}
\title{Public Request Contracts for Source-Only Releases\\\large Exact Paper-Level Handoffs from Boilerplate Policy to Maintained Packet Cuts}
\author{Anonymous}
\date{Unified archive note v0.10 (2026-03-21; archive v466)}
\begin{document}
\maketitle

\begin{abstract}
A short paper in this archive often says only one public sentence about omitted support: the release is source-only and supporting artifacts are available on request. That sentence is useful but still leaves one repeat question for referees and later maintainers: exactly which maintained packet/menu/validator cut stands behind that public promise for this paper? This note gives that bridge one home. It defines a compact \emph{public request contract}: one paper-level object that binds archive-wide source-only boilerplate to the already-maintained disclosure packets, review menus, validator companions, requestable evidence classes, and public anchor families that make the sentence auditable rather than aspirational.
\end{abstract}

\paragraph{Non-redundancy note.}
Synthesis~6 owns the archive-wide source-only artifact policy, Synthesis~52 owns exact answer disclosure packets, Synthesis~54 owns answer review menus, Synthesis~55 owns the theorem-to-review bridge, Synthesis~57 owns reusable omitted-support family names, Synthesis~58 owns exact completion criteria for satisfying such promises, and Synthesis~17 owns the concrete worked instantiation.\cite{series:artifactpolicy,series:challengeanswerdisclosurepackets,series:challengeanswerreviewmenus,series:seriesspines,series:requestableevidenceclasses,series:requestfulfillmentcertificates,series:workedexample}
This note owns only the paper-facing aggregation that says which maintained packet/menu/validator slice a particular source-only public sentence is actually promising. Stable omitted-support family names for those slices belong in Synthesis~57 rather than here, exact completion criteria for when a promised cut has been fully satisfied belong in Synthesis~58 rather than here, the within-release paper-level public status token for the overall promise belongs in Synthesis~59 rather than here, cross-release carryforward of that public token belongs in Synthesis~60 rather than here, the exact outward-facing notice sentence readers should see belongs in Synthesis~61 rather than here, the canonical sentence family plus slot realization rule behind that wrapper belong in Synthesis~62 rather than here, the token-keyed choice of which such family applies belongs in Synthesis~63 rather than here, the request-side terminal stop belongs in Synthesis~73 rather than here, the paired answer/request citation discipline belongs in Synthesis~74 rather than here, and the stop-versus-widen/reopen rule belongs in Synthesis~75 rather than here.

\paragraph{Scope.}
This is not a new disclosure packet, review menu, escalation ladder, or release stage.
It is one paper-level handoff contract saying how a public ``supporting artifacts are available on request'' sentence resolves into already-maintained objects for one paper.

\paragraph{Exports.}
This note exports (i) public request contracts, (ii) a policy-anchor rule, (iii) a public-anchor-family rule, (iv) a shared-public-anchor-family-vocabulary rule, and (v) a maintained-cut rule saying that later papers may cite one compact object when they need an exact, auditable meaning for their source-only public request sentence.

\paragraph{Later-terse-paper citation posture.}
Later terse papers should cite Synthesis~56 only when the live issue is the exact paper-facing source-only promise itself: which public-anchor families one paper actually promises on request, which maintained disclosure/review slices and validator companions that promise binds, or which compact contract key later route and packet mirrors are merely importing. Stop at Synthesis~57 when the live issue is the reusable omitted-support family typing behind that promise. Widen to Synthesis~58 when the live issue is whether one promised family cut has actually been fulfilled, widen to Synthesis~59--60 when the live issue is the within-release paper-level status token or its cross-release carryforward, stop at Synthesis~61--63 when the contract basis is already fixed and the remaining issue is the exact outward-facing sentence, canonical notice family, or token-keyed family selector shown to readers, and stop at Synthesis~17 when the concrete maintained worked route is the live issue instead. Under the paired terminal citation discipline and paired cutover rule, later terse papers should usually stop at Synthesis~73--75 once the live question has become the request-side terminal stop, the paired reuse rule, or the stop-versus-widen judgment rather than reopen this paper-facing contract note.\cite{series:requestableevidenceclasses,series:requestfulfillmentcertificates,series:publicrequeststatusenvelopes,series:publicrequestcarryforward,series:publicrequestnotices,series:publicrequestnoticenormalforms,series:publicrequestnoticeselections,series:requestsideterminalcitationnormalform,series:pairedterminalcitationdiscipline,series:pairedterminalcutoverrules,series:workedexample}

\section{Why public request contracts deserve their own note}
The archive-wide artifact policy intentionally stays broad.
Disclosure packets intentionally stay answer-local.
Review menus intentionally stay request-class-local.
Series spines intentionally stay bridge-like.
A short source-only paper often needs all of those, but only in one tiny paper-facing combination.
Without one maintained paper-level contract, that combination tends to drift into prose: one paper says ``ask for the archive,'' another says ``ask for the validator bundle,'' and a third says nothing beyond generic boilerplate.
That is exactly the kind of repetition and ambiguity this series is trying to remove.

\begin{definition}[Public request contract]
Fix one paper key $p$ in a source-only public release.
A \emph{public request contract}
\[
R(p)=(p,A,F,M,V,K,\rho)
\]
contains an archive-wide policy anchor $A$, a public-anchor-family set $F$, a maintained-cut set $M$ naming the already-maintained disclosure packets and review menus that back the paper's public sentence, a validation-companion set $V$, one evidence-class-key family $K$, and one canonical request phrase $\rho$.
The contract does not create new answer objects.
It only states the exact maintained meaning of the paper's public request promise; reusable omitted-support family names remain imported evidence classes rather than new packet ownership.
\end{definition}

\begin{definition}[Policy-anchor rule]
A public request contract satisfies the \emph{policy-anchor rule} if its policy anchor names the archive-wide source-only artifact policy and does not restate or override that policy.
So the paper-level contract narrows the meaning of the sentence for one paper without duplicating archive-wide boilerplate ownership.
\end{definition}

\begin{definition}[Public-anchor-family rule]
A public request contract satisfies the \emph{public-anchor-family rule} if every public anchor family it lists resolves to a public path that is actually cited or shipped for that paper and to the corresponding disclosure-packet and review-menu entries that govern follow-up on that same family.
So the contract names not just hidden support paths, but also the public anchors by which a reader first encounters the claim.
\end{definition}

\begin{definition}[Shared-public-anchor-family-vocabulary rule]
A public request contract satisfies the \emph{shared-public-anchor-family-vocabulary rule} if every public-anchor-family key it defines may recur in later evidence-class ledgers, fulfillment certificates, routes, or series-spine mirrors only as imported contract vocabulary: those later objects may point back to one stable family basis, but they do not thereby re-own the public path set, review-menu binding, or disclosure-packet binding attached to that family.
So repeated family keys stay reusable without silently moving public-anchor ownership out of the contract layer.
\end{definition}

\begin{definition}[Maintained-cut rule]
A public request contract satisfies the \emph{maintained-cut rule} if every listed review menu, disclosure packet, support-manifest path, artifact-inventory path, and validator companion already exists as a maintained object in the same worked release.
So the contract is an auditable aggregation of maintained objects, not a prose promise to assemble them later.
\end{definition}

\begin{lemma}[Public request contracts stabilize source-only handoffs without widening ownership]
If a paper cites an archive-wide artifact policy and a public request contract satisfying the policy-anchor, public-anchor-family, shared-public-anchor-family-vocabulary, and maintained-cut rules, then the paper's public source-only sentence has one exact auditable interpretation: a reader can identify the public anchor families involved, the smallest maintained packet/menu objects that govern the relevant follow-up classes, the evidence-class names for the omitted-support families if the contract imports them, and the validator/manifests that fix the maintained bundle cut.
\end{lemma}

\begin{table}[t]
\centering
\small
\begin{tabularx}{\linewidth}{@{}l l X@{}}
\toprule
Question family & First sufficient note & Widen only if \\
\midrule
What exactly did this paper promise on request? & Synthesis~56 & The question becomes reusable family typing, completion, paper-level status, or visible notice choice rather than the paper-facing contract itself. \\
Which stable omitted-support family names are in play? & Synthesis~57 & The question becomes whether those families are complete for this release cut. \\
Is that promised cut fully satisfied? & Synthesis~58 & The question becomes compression of completion into a paper-level public token. \\
What is the current paper-level status, and does it carry forward? & Synthesis~59--60 & The question becomes the visible sentence / canonical family chosen for that status basis. \\
What exact outward-facing notice sentence/family was shown? & Synthesis~61--63 & The question fixes a follow-up class and becomes packetization rather than visible notice choice. \\
\bottomrule
\end{tabularx}
\caption{The first source-only owner bucket after the theorem-to-review bridge can be walked without prose drift by stopping at the first note that answers the current question family.}
\label{tab:publicrequestcontractstoprules}
\end{table}

\section{A minimal public request-contract card}
\begin{table}[t]
\centering
\small
\begin{tabularx}{0.97\linewidth}{@{}p{0.31\linewidth}X@{}}
\toprule
\textbf{Public row} & \textbf{Pinned value}\\
\midrule
Policy anchor & \texttt{policy\_anchor.cite\_key = series:artifactpolicy}; later terse papers should quote the contract's imported policy anchor rather than restate source-only boilerplate here \\
Worked contract floor & \texttt{public\_request\_contract\_key = worked\_example\_receipt\_interlock\_public\_request\_contract} from \nolinkurl{example_public_request_contracts.json} \\
Request phrase floor & ``This paper is source-only; supporting artifacts are available on request.'' \\
Public-anchor-family floor & \texttt{public\_status\_sentence\_family} and \texttt{compact\_diff\_summary\_family} with first public anchors in \nolinkurl{example_successor_lineage_notice.json}, \nolinkurl{example_successor_claim_support_map.json}, \nolinkurl{example_compare_report.json}, and \nolinkurl{example_successor_clause_pack.json} \\
Maintained disclosure / review floor & The contract binds the four worked disclosure packets and four worked review menus rather than leaving those paper-facing handoffs implicit in prose \\
Imported evidence-class floor & The contract imports \texttt{bundle\_identity\_and\_validator\_companions}, \texttt{public\_status\_lineage\_and\_closure\_support}, and \texttt{compact\_diff\_compare\_and\_replay\_support} as stable omitted-support names rather than re-enumerating support paths inline \\
Completion / validation floor & The contract points forward to \texttt{public\_status\_sentence\_request\_fulfillment\_certificate} and \texttt{compact\_diff\_summary\_request\_fulfillment\_certificate} and carries \path{support_manifest.json}, \path{example_artifact_inventory.json}, \path{example_validation_report.json}, and \path{../tools/validate_example.py} as its validation companions \\
Escalation pointer & Widen to Synthesis~57 for family typing, Synthesis~58 for family completion, Synthesis~59--63 for paper-level status / notice layers, and Synthesis~73--75 once the live issue is terminal reuse or cutover rather than the paper-facing promise \\
\bottomrule
\end{tabularx}
\caption{A minimal public request-contract card. This is the smallest public answer later terse papers can quote directly when the live issue is what one source-only paper actually promised on request and which maintained families, packets, menus, and validator companions that promise binds.}
\label{tab:minimal-public-request-contract-card}
\end{table}

This card is intentionally non-decorative: it pins one worked contract key, the policy anchor it imports, the exact request phrase, the two public-anchor families the promise covers, the imported evidence classes, and the validator-companion floor without silently re-owning family typing, family completion, paper-level status compression, or request-side terminal judgment.\cite{series:workedexample}

\paragraph{One tiny request-contract drill.}
In the worked note, the same contract key covers two different public-anchor families at once: the public-status lineage sentence and the compact-diff summary sentence. That single paper-facing promise is honest only because the contract keeps both family keys visible, keeps the shared bundle-identity / validator class visible, and still points forward to the family-scoped completion certificates rather than pretending that the contract itself proves fulfillment. So a later terse paper may quote this contract card directly when the live issue is the exact paper-facing promise. If it instead needs to know whether the compact-diff family is already complete, it should widen immediately to Synthesis~58; if it needs the exact visible sentence readers were shown, it should widen to Synthesis~61--63 rather than stretching the contract into a notice note.\cite{series:workedexample}

\section{Worked example pointer}
The maintained worked bundle now ships \nolinkurl{example_public_request_contracts.json}.\cite{series:workedexample}
It records one paper-facing contract for the worked note, tying the archive-wide source-only policy anchor to the public status and compact-diff anchor families, their answer disclosure packets, their answer review menus, the requestable evidence classes that name the omitted-support families, the companion request-fulfillment certificates that say when those promised cuts have been fully satisfied, and the validator companions that make the on-request promise auditable. Each listed public anchor family now also exports an explicit vocabulary-posture tag, and the worked route/spine surfaces mirror that ledger through a shared \nolinkurl{source_only_public_anchor_family_mode_map} rather than silently flattening family basis into later evidence, completion, or status objects.\cite{series:requestfulfillmentcertificates}
So a later short paper can cite one maintained paper-facing object instead of restating boilerplate policy, packet paths, menu choices, completion criteria, and validator hooks by hand. The worked route catalog now imports that same contract key beside each shipped audience/question route, so a terse downstream note can recover the paper-level omitted-support promise from the same maintained object that already names its checked sentence and referee-side handoff surfaces without pretending that the route owns the contract or the public-anchor family basis it merely mirrors. Under the paired terminal citation discipline and paired cutover rule, later terse papers should usually stop at Synthesis~73--75 for the route-side stop/widen judgment and reopen this contract note only when the exact paper-facing promise, anchor-family basis, or maintained-cut bundle identity is itself disputed.\cite{series:requestsideterminalcitationnormalform,series:pairedterminalcitationdiscipline,series:pairedterminalcutoverrules}

\begin{thebibliography}{99}\small

\bibitem{series:artifactpolicy}
Anonymous.
\newblock Artifact and Disclosure Policy for the One-True-Archive (Source-Only Public Release).
\newblock Series synthesis note (Synthesis~6), 2026. (This archive.)

\bibitem{series:challengeanswerdisclosurepackets}
Anonymous.
\newblock Challenge Answer Disclosure Packets for Successor Maintenance: Public Pointers and Requestable Audit Slices under Space Pressure.
\newblock Series synthesis note (Synthesis~52), 2026. (This archive.)

\bibitem{series:challengeanswerreviewmenus}
Anonymous.
\newblock Challenge Answer Review Menus for Successor Maintenance: Smallest-Sufficient Referee Handoffs from Dockets, Profiles, Packets, and Ladders.
\newblock Series synthesis note (Synthesis~54), 2026. (This archive.)

\bibitem{series:seriesspines}
Anonymous.
\newblock Series Spines from Mathematical Roots to Referee Handoffs: A Non-Redundant Bridge from Backbones, Receipts, Replay, Release Evolution, and Review Menus.
\newblock Series synthesis note (Synthesis~55), 2026. (This archive.)

\bibitem{series:requestableevidenceclasses}
Anonymous.
\newblock Requestable Evidence Classes for Source-Only Releases: Stable Names for Omitted-Support Families behind Packet Cuts and Public Request Contracts.
\newblock Series synthesis note (Synthesis~57), 2026. (This archive.)

\bibitem{series:requestfulfillmentcertificates}
Anonymous.
\newblock Request Fulfillment Certificates for Source-Only Releases: Exact Family-Completion Witnesses behind Public Request Promises.
\newblock Series synthesis note (Synthesis~58), 2026. (This archive.)

\bibitem{series:publicrequeststatusenvelopes}
Anonymous.
\newblock Public Request Status Envelopes for Source-Only Releases: Paper-Level Tokens Compressing Family Completion without Re-owning the Promise.
\newblock Series synthesis note (Synthesis~59), 2026. (This archive.)

\bibitem{series:publicrequestcarryforward}
Anonymous.
\newblock Public Request Carryforward for Source-Only Successor Releases: When Paper-Level Request Status Persists, Advances, or Reopens across Release Evolution.
\newblock Series synthesis note (Synthesis~60), 2026. (This archive.)

\bibitem{series:publicrequestnotices}
Anonymous.
\newblock Public Request Notices for Source-Only Releases: Outward-Facing Sentences for Within-Release Status and Cross-Release Carryforward.
\newblock Series synthesis note (Synthesis~61), 2026. (This archive.)

\bibitem{series:publicrequestnoticenormalforms}
Anonymous.
\newblock Public Request Notice Normal Forms for Source-Only Releases: Canonical Sentence Families for Within-Release Status and Cross-Release Carryforward Notices.
\newblock Series synthesis note (Synthesis~62), 2026. (This archive.)

\bibitem{series:publicrequestnoticeselections}
Anonymous.
\newblock Public Request Notice Selections for Source-Only Releases: Choosing Canonical Notice Families from Within-Release Status and Cross-Release Carryforward Tokens.
\newblock Series synthesis note (Synthesis~63), 2026. (This archive.)

\bibitem{series:requestsideterminalcitationnormalform}
Anonymous.
\newblock Public Request Response Packet Refresh Response Packet Closure Verdicts for Source-Only Successor Releases: Terminal Closure Stops for the Visible Refresh Branch.
\newblock Series synthesis note (Synthesis~73), 2026. (This archive.)

\bibitem{series:pairedterminalcitationdiscipline}
Anonymous.
\newblock Paired Terminal Citation Discipline for Stable Review Spines: One Answer-Side Stop and One Request-Side Capstone for Terse Papers.
\newblock Series synthesis note (Synthesis~74), 2026. (This archive.)

\bibitem{series:pairedterminalcutoverrules}
Anonymous.
\newblock Paired Terminal Cutover Rules for Stable Review Spines: When to Stop at the Checked Answer, When to Widen to Source-Only, and When to Fail Closed to the Reopening Owner.
\newblock Series synthesis note (Synthesis~75), 2026. (This archive.)

\bibitem{series:workedexample}
Anonymous.
\newblock Worked Example: Receipt Interlock for an Anonymous-DHT Reader-Privacy Claim.
\newblock Series synthesis note (Synthesis~17), 2026. (This archive.)

\end{thebibliography}
\end{document}
